\section*{Overview}

Lazy propagation is what turns a plain segment tree into a practical range-update structure. Instead of pushing every
update all the way to the leaves immediately, the tree delays work until that work is actually needed.

This is the standard contest upgrade when the problem asks for range updates together with range queries, such as
``add \(x\) to every value on \([l, r]\)'' and ``report the sum on \([l, r]\)''.

\section*{When to Use It}

Use it when:

\begin{itemize}
  \item updates affect whole intervals rather than single points,
  \item queries still ask for segment aggregates,
  \item the update can be composed compactly as a lazy tag,
  \item doing the update naively would touch too many leaves.
\end{itemize}

If the operation is static, a difference array might be enough. If the update logic is too complicated to summarize in
one tag, a plain lazy tree may stop being the right tool.

\section*{Core Idea}

Each node still represents a segment and stores its aggregate. The new ingredient is a \textbf{lazy tag}: pending work
that conceptually applies to the whole segment but has not yet been pushed to the children.

For range add / range sum, if I add \(x\) to a segment of length \(\ell\), then:

\begin{itemize}
  \item the node sum increases by \(x \cdot \ell\),
  \item the children do not need to be updated immediately,
  \item I only record that both children should eventually receive \(x\) as well.
\end{itemize}

That deferred record is the lazy value.

\section*{Key Insight}

Lazy propagation works because a query or update never needs the children to be exact \emph{until} it descends below
the current node. As long as the whole current segment is either fully covered or fully skipped, it is enough that the
current node summary is correct.

So the tree can postpone detailed child work without losing correctness. The delayed work is pushed only when a later
operation partially overlaps the segment and needs to see finer structure.

\section*{Operations / Main Technique}

The three core routines are:

\begin{itemize}
  \item \textbf{apply(node, tag)}: update the node summary and combine the lazy tag,
  \item \textbf{push(node)}: send the pending tag to the children before going deeper,
  \item \textbf{pull(node)}: recompute the parent from the children after a recursive change.
\end{itemize}

\paragraph{Worked example.}
Suppose the tree stores sums and I apply ``add \(5\)'' on \([1, 8]\). The root sum increases by \(5 \cdot 8\), and the
root lazy tag becomes \(+5\). If the next query asks only about \([1, 4]\), then the root must first \emph{push} the
tag so the left and right children both learn about that pending \(+5\).

\section*{Correctness Intuition}

The invariant is that every node's stored summary already reflects all updates that fully covered that node's segment,
even if some of those updates have not been propagated to descendants yet.

When we descend, \texttt{push} preserves that invariant by transferring the pending effect to the children. When we
return from recursion, \texttt{pull} restores the parent summary from the now-correct child summaries.

So lazy propagation does not skip work forever. It only postpones work until the tree actually needs lower-level
detail.

\section*{Complexity Analysis}

\begin{itemize}
  \item range update: \(O(\log n)\),
  \item range query: \(O(\log n)\),
  \item memory: \(O(n)\).
\end{itemize}

Each operation visits only \(O(\log n)\) nodes on the recursion frontier, and each visited node does only constant
extra work for tag handling.

\section*{Implementation}

The reference code implements the most common teaching version:

\begin{itemize}
  \item range add,
  \item range sum,
  \item recursive tree,
  \item one lazy tag per node.
\end{itemize}

That version is useful because the meaning of every field is easy to state. Once this shape is stable, it is much
easier to adapt it to range assign, range min/max, or tree-path problems.

\section*{Common Pitfalls}

\begin{itemize}
  \item Updating the node sum but forgetting to store the lazy tag.
  \item Pushing too late, so a partial-overlap recursion reads stale child data.
  \item Pulling from children that have not yet received the pushed tag.
  \item Mixing segment length formulas, especially with inclusive endpoints.
  \item Treating range assign and range add as if they compose the same way.
\end{itemize}

Most wrong answers come from violating the meaning of ``what is already reflected in this node''.

\section*{Variants / Extensions}

\begin{itemize}
  \item range assign + range sum,
  \item range min/max with range add,
  \item lazy segment tree over custom monoids,
  \item tree-path queries after heavy-light decomposition,
  \item segment tree beats when one lazy tag is no longer expressive enough.
\end{itemize}

AtCoder Library's \texttt{lazysegtree} is a good reference for the ``generic algebraic'' viewpoint, but in contest
practice I still like understanding one concrete hand-written version first.

\section*{Practice Problems}

\begin{itemize}
  \item Range add and range sum on one array.
  \item Range assign and range minimum.
  \item Heavy-light decomposition problems that need path updates instead of only path queries.
  \item Problems where you must define the node summary and lazy tag yourself.
\end{itemize}

\section*{References}

\begin{itemize}
  \item \href{https://cp-algorithms.com/data_structures/segment_tree.html#range-updates-lazy-propagation}{cp-algorithms: Lazy propagation}.
  \item \href{https://atcoder.github.io/ac-library/production/document_en/lazysegtree.html}{AtCoder Library: lazysegtree}.
  \item \href{https://github.com/kth-competitive-programming/kactl}{KACTL}.
  \item \href{https://codeforces.com/catalog}{Codeforces Catalog}.
\end{itemize}
